\subsection{Good models for already projective relative adjoints}

\begin{proposition}[Projective relative completion]
\label{prog:projective-good-model}
Let \(f:X\to S\) be an already projective surjective morphism of
normal compact Kähler spaces, with \(X\) smooth. Let
\((X,B+\mathbf M)\) be an effective generalized klt pair with adjoint
\(A\), carried by a projective resolution \(p:W\to X\) with smooth
\(W\), so that
\[
 p^*A=c_1(K_W)+[B_W]+\mathbf M_W.
\]
Suppose that there are an actual real line bundle
\(J\in\Pic(X)\otimes\R\) and
\(\beta\in H^{1,1}_{\BC}(S)\) such that
\[
 A=c_1(J)+f^*\beta,
 \qquad
 N:=p^*J-K_W-B_W
 \text{ is nef and big over }fp.
\]
Assume that \(J\) is pseudoeffective over \(S\).
Then there is a finite chosen \(A\)-negative projective program over
\(S\), with globally strongly \(\Q\)-factorial compact Kähler
working spaces, and a nonextracting endpoint \(X_m\) such that
\[
 J_m\sim_{\R}g^*H,
 \qquad
 A_m=g^*\bigl(c_1(H)+\rho^*\beta\bigr).
\]
Here \(J_m\) is the actual real-line transform of \(J\),
\(g:X_m\to Z\) is a projective connected-fibre morphism,
\(\rho:Z\to S\) is projective, and \(H\) is an actual real line
bundle ample over \(S\). Thus the class on \(Z\) in this formula is
relatively Kähler over \(S\). The generalized pair on \(X_m\) is
gklt. On a common resolution the comparison with \(A\) is effective
and exceptional over \(X_m\), with strictly positive coefficient
at the strict transform of every prime contracted from \(X\).
\end{proposition}

\begin{proof}
We reduce to fixed ordinary pairs, and construct one global program.
Choose finitely many Stein open sets \(U_\ell\subset S\) and
semianalytic Stein compact subsets \(W_\ell\Subset U_\ell\) whose
interiors cover \(S\). Such compacts are obtained by intersections
with closed polydiscs in local embeddings of \(S\). In particular
they satisfy condition~(P) of \cite{FujinoAnalyticBCHM}: the source
is normal, the base is Stein, the compact is Stein, and its
intersection with any analytic subset defined near it has finitely
many connected components. This remains true on every normal
projective model. Shrinking \(U_\ell\) around \(W_\ell\) preserves
the finite cover.

\smallskip\noindent
\emph{Actual ordinary replacements.}
On each chart represent the finitely many line bundles in the data
by Cartier divisors. For a projective map to a Stein space this can
be done by twisting a line bundle and its prospective meromorphic
frame by sufficiently positive relative lines and using nonzero
sections. The generic-rank-one direct-image argument is an
alternative only when the map is bimeromorphic.
Relative bigness gives
\(N\sim_{\R}P+E\), with \(P\) relatively ample and \(E\geq0\).
For sufficiently small \(\epsilon>0\),
\((W,B_W+\epsilon E)\) is sub-klt near the compact inverse image.
The line \((1-\epsilon)N+\epsilon P\) is relatively ample.
General divided relative sections give an effective representative
\(T\) such that
\[
 \Theta:=\epsilon E+T\sim_{\R}N,
 \qquad (W,B_W+\Theta)\text{ is sub-klt}.
\]
All supports and coefficient margins are fixed near that compact.
With compatible canonical representatives, push the finite
principal-divisor expression for \(\Theta-N\) down and pull it
back again. For \(\Delta_\ell=p_*(B_W+\Theta)\) this gives
\[
 \Delta_\ell\geq0,
 \qquad K_X+\Delta_\ell\sim_{\R}J|_{X_{U_\ell}},
 \qquad
 p^*(K_X+\Delta_\ell)=K_W+B_W+\Theta.
\]
In the last equality the chosen principal correction is included
in the representative of \(K_X+\Delta_\ell\).
Thus \((X_{U_\ell},\Delta_\ell)\) is ordinary klt and its boundary
is relatively big. Lemma~11.15 of \cite{FujinoAnalyticBCHM}, with
its big part equal to \(\Delta_\ell\) and its remaining part zero,
replaces it, in its actual real linear equivalence class, by a klt
boundary with an effective general relatively ample rational
summand. Its condition on non-klt centres is vacuous.
Choose the reserved ample summand to be \(\Q\)-linearly equivalent
to a small rational multiple of a fixed global \(f\)-ample line.
Its later strict transform is consequently \(\Q\)-Cartier by the
transport of that global line, without a local factoriality claim.

We make these choices for a finite rational polytope of actual
global line data. Start with a finite rational span containing
\(J\), and add relatively ample rational lines spanning the global
relative degree space. The above construction is valid in a
neighbourhood of \(J\): retain a small part of the ample summand
and use openness of ampleness and the strict klt inequalities.
It is also valid along a segment from \(J\) to a sufficiently
positive adjoint, by adding divided general ample divisors.
Use finitely many such representatives and one log resolution on
each chart. More explicitly, represent a finite simplex inside the
open ample neighbourhood used above; varying its positive
coefficients represents all nearby global line parameters on the
same finite support. Form the joint finite affine system consisting
of the divisor identities on all charts, with common global line
coordinates and separate boundary and principal-divisor
coordinates. This system is rational. Allow all boundary
coefficients to vary, including those originally contributed by
the real boundary \(B\); fix only the reserved rational ample part
and the zero-support constraints. Its projection to the global
line coordinates contains the just constructed neighbourhood, so
it has a rational affine section. Choose that section sufficiently
close to the constructed real affine lift, by rational approximation
in its affine space of sections, so that the strict klt and
effectiveness margins are retained along the compact segment.
In the joint solution space the
strict klt and effectiveness margins permit a rational polytope
containing the scaling segment. This supplies on every chart
a rational polytope of ordinary klt boundaries with a fixed
effective relatively ample rational summand, representing the
restrictions of the same global line parameters. These boundaries,
including the scaling representatives, are fixed before the
program starts. Their divisors need not agree on different charts.

\smallskip\noindent
\emph{One global scaling construction.}
Apply the global degree-space construction in the proof of
Proposition~\ref{prog:relative-klt} to \(J\) and a general ample
scaling direction in the chosen span. Here are the modifications
that allow real line data and a nonbirational map to \(S\).
The global relatively ample line still gives a compact normalized
slice of the curve cone. The fixed ordinary chart pairs give finite
negative-ray truncations. A general direction avoids equal wall
parameters for independent ray-degree functionals, exactly as in
that proof. A positive wall therefore selects a single ray in the
global degree space. A rational nef support is represented by a
global rational line; a large multiple minus the chartwise ordinary
adjoint is relatively ample. Ordinary relative base point freeness
contracts the ray using the evaluation of this one global line.

If the contraction is small, choose a global rational line having
negative ray degree and construct its flip by its one global
graded algebra. The local ordinary proof of
Proposition~\ref{prog:detected-flip} applies to this algebra:
all curves of the contraction have proportional degrees for global
lines, so the detector and the driving line are positively
proportional over the contraction; ordinary rational replacement
gives local finite generation. This part of that proof needs only
the already projective contraction and these line degrees.
It does not require a prior assertion about the full Bott--Chern
cone. The relative Proj is the next global model. In a divisorial
step the exceptional prime and
Lemma~\ref{prog:integral-descent} give the same continuation as in
Proposition~\ref{prog:relative-klt}. Killing the degree of any global
line by a rational multiple of the detector, then applying that
lemma, preserves global strong \(\Q\)-factoriality in both cases.
The global lines continue to span the relative degree spaces.
The rational support-simplex argument in that proposition makes
the wall a descended relatively ample real line on its contraction
base and gives a nonempty interval of ampleness immediately after
each step. This argument uses relative line data and projectivity,
not birationality of the map to \(S\).

We check explicitly the persistence of the ordinary pairs. Expand
the line parameters in finitely many global rational lines. Their
reflexive transforms are actual rational lines by the preceding
strong property. Pushing forward the fixed principal-divisor
identities gives
\[
 K_{X_i}+\Delta_{\ell,i}\sim_{\R}J_i|_{(X_i)_{U_\ell}}.
\]
Hence these ordinary adjoints are \(\R\)-Cartier; no local
factoriality of individual prime divisors is being assumed.
On a common resolution their comparison minus the appropriate
positive multiple of the detector comparison is exceptional and
relatively numerically trivial. Negativity in both signs makes
this difference zero. Every chosen step is consequently negative
for the fixed ordinary pairs. They remain effective and klt, and
their boundaries remain relatively big. The same equality of
comparison divisors, after adding the fixed base pullback,
preserves the generalized pair and its discrepancy inequalities.

For termination, every already constructed working model is a weak
log canonical model on each chart of a boundary in the fixed
polytope: choose a parameter in its nonempty ample scaling interval.
Indeed, actual descent at a preceding wall \(\lambda_j\) makes
the comparison for \(J+tH\) equal to
\((1-t/\lambda_j)\) times its \(J\)-comparison. For a parameter
in the current interval all preceding \(\lambda_j\) are at least
\(t\), so these comparisons are effective.
Theorem~E of \cite{FujinoAnalyticBCHM} gives finitely many marked
models on each chart. That theorem does not require locally
\(\Q\)-factorial targets. There are therefore finitely many possible
global marked working models. Indeed, isomorphisms between the
restrictions of two existing marked models agree on their common
dense open, and hence on overlaps. A repeated marked working model
would have the same boundary transform and discrepancies, contrary
to strict discrepancy increase at an intervening nontrivial step.
The global program is finite. Relative pseudoeffectivity excludes
a Mori fibre endpoint, so \(J_m\) is nef over \(S\).
This is a finite-cover argument for a single constructed program;
it does not glue independent minimal models.

\smallskip\noindent
\emph{Semiampleness and the global polarization.}
On the endpoint the fixed ordinary pair is klt, its effective
boundary is relatively big, and its adjoint is nef over
\(W_\ell\). Lemma~11.16 of \cite{FujinoAnalyticBCHM} applies to
this already existing weak model. More explicitly, Lemma~11.15
replaces the endpoint boundary by \(A_\ell+B_\ell\), where
\(A_\ell\geq0\) is relatively ample, \(B_\ell\geq0\), and the
pair is klt. Theorem~8.3 then gives semiampleness near \(W_\ell\):
its log canonical, klt-existence, real Cartier, ample-summand and
nef-over-compact hypotheses have all been checked.

The actual equivalences identify these local semiample adjoints
with restrictions of \(J_m\). Their ample models agree on overlaps
by \cite[Lemma~11.3]{FujinoAnalyticBCHM}. To retain a global
polarization, shrink the initial parameter polytope around \(J\)
after the finite program: all its finitely many strict step signs
persist, so the transformed ordinary boundaries remain klt and big.
On each chart the pushforward of its initial common ample rational
summand is effective and relatively big. Apply
\cite[Lemma~11.13]{FujinoAnalyticBCHM} to this common big rational
part and the small endpoint boundary polytope. All pairs are klt,
so the condition on non-klt centres is vacuous. The lemma gives a
common general relatively ample rational summand by one
\(\Q\)-linearly trivial translation. In particular the rational
parametrization is preserved; separate applications of
Lemma~11.15 to individual parameters would not suffice here.
In this finite rational span Theorem~11.17 of \cite{FujinoAnalyticBCHM}
makes the endpoint nef region a rational polytope. Intersect the
finitely many inverse images of these regions and express \(J_m\)
as a positive combination of rational points in its minimal face.
The corresponding global rational lines are semiample on every
chart by the same ordinary base point free argument. After clearing
denominators and taking common sufficiently divisible powers over
the finite cover, each is relatively generated globally: generation on
the charts is generation of its global coherent relative
evaluation map. The product of these finitely many relative
morphisms, followed by Stein factorization, gives a projective
\(g:X_m\to Z\), a projective \(\rho:Z\to S\), and descended
global lines whose positive combination \(H\) is relatively ample.
It realizes the glued ample model and gives
\(J_m\sim_{\R}g^*H\). Adding back \(f_m^*\beta\) proves the
asserted Bott--Chern identity. Composing the effective step
comparisons proves the final discrepancy statement.
\end{proof}

\begin{corollary}[Semiampleness on an existing projective model]
\label{prog:projective-semiample}
Let \(f:T\to S\) be a projective morphism of normal compact Kähler
spaces with \(T\) globally strongly \(\Q\)-factorial, and let
\(J\in\Pic(T)\otimes\R\) be nef over \(S\). Suppose that on a
finite Stein-compact cover satisfying condition~(P), there are
effective ordinary klt boundaries \(\Delta_\ell\), big over the
base, such that
\(K_T+\Delta_\ell\sim_{\R}J|_{T_{U_\ell}}\).
Then \(J\) has the projective relative semiample contraction and
actual descended ample real line of
Proposition~\ref{prog:projective-good-model}.
In particular this applies to any already existing weak model of
the fixed ordinary chart pairs in that proof; one need not run a
new program for the adjoint in question.
\end{corollary}

\begin{proof}
Lemma~11.15 and Theorem~8.3 of \cite{FujinoAnalyticBCHM} give
local semiampleness. For an existing weak model this is also
Lemma~11.16. To globalize the polarization, take a finite rational
span of the global lines defining \(J\). On each chart the ample
part of the normalized boundary allows small perturbations in
this span. The joint rational affine construction in the preceding
proof gives a rational family with a common rational ample part;
all boundary coefficients are allowed to vary. Theorem~11.17 then
gives a rational nef region on each chart. The finite-intersection,
rational-vertex and global-evaluation argument in that proof gives
one projective contraction and its actual descended ample real
line. Only its local ample models are identified by uniqueness.
\end{proof}

\begin{remark}[The two uses of projective completion]
\label{prog:projective-good-model-uses}
For a selective extraction over an already constructed weak model,
the carrier map is projective bimeromorphic. The inherited
exceptional splitting expresses the nef datum, modulo this base,
by actual real lines. Relative nefness follows from the carrier
datum, while relative bigness follows from birationality, using the
ample-plus-effective construction in the ordinary replacement.
The prepared adjoint is a base pullback plus an effective
exceptional divisor supported on exactly the unwanted primes.
Proposition~\ref{prog:projective-good-model} applies. On its endpoint
the remaining error is effective, exceptional over the fixed base,
and relatively nef; negativity makes it zero. Strict discrepancy
improvement prevents contraction of a requested zero-error prime.
The chosen extraction boundary must be effective and klt; in
particular its prescribed primes have the usual admissible
discrepancy range.

For the graph construction, suppose the already projective graph
resolution \(h:V\to Y_m\) has smooth source and
\[
 A_V=h^*A_{Y_m}+[G],\qquad G\geq0.
\]
Then the global relative real line is \(J=\OO_V(G)\), the base
class is \(A_{Y_m}\), and its carrier line is
\(N=G-K_V-B_V\). Its relative degrees equal those of the prepared
nef datum. That datum is nef and big on the prepared carrier, for
example a pullback of its prepared Kähler class; thus \(N\) is nef
and big over \(Y_m\). This bigness is a separate hypothesis check,
since \(h\) can have positive-dimensional general fibres.
Effectivity of \(G\) gives relative pseudoeffectivity. Apply the
proposition over \(Y_m\). The subsequent negativity argument must
still remove \(G_m\); only then is the endpoint adjoint the pullback
of \(A_{Y_m}\) and good over the original base.

For transport of all classes, verify separately that all
Bott--Chern classes on the graph carrier are real-line classes
modulo \(Y_m\), using the smooth rationally connected fibration
and the inherited exceptional splitting on the lower-dimensional
base model. Lemma~\ref{prog:projective-cone} then identifies
the relative degree rays with full Bott--Chern rays, so the
detected-step transport applies. The driving identity involving
\(G\) alone does not give this additional property.
\end{remark}

\begin{corollary}[Ordinary projective Mori output]
\label{prog:projective-mori}
Let \(f:X\to S\) be an already projective surjective morphism of
compact Kähler spaces, with \(X\) smooth. Let \((X,B)\) be an
effective ordinary real klt pair, and let \(D=K_X+B\). If \(D\)
is not pseudoeffective over \(S\), then its projective program over
\(S\), with scaling of a relatively ample actual real line \(N\),
terminates with a Mori fibre contraction. The initial scaling
parameter is chosen so that \(D+TN\) is nef over \(S\).
The working spaces are globally strongly \(\Q\)-factorial and
compact Kähler, and the actual line comparisons and integral
descent of Proposition~\ref{prog:projective-good-model} apply.
\end{corollary}

\begin{proof}
The line \(N\) is fixed; it need not be a general scaling direction.
In particular several extremal rays may have the same wall
parameter. Choose \(\epsilon>0\) such that \(D+2\epsilon N\) is
still not pseudoeffective over \(S\). On the finite Stein-compact
cover choose divided general effective representatives of \(N\)
so that all the resulting ordinary pairs representing \(D+tN\),
\(0\leq t\leq T\), are klt. For
\(\epsilon\leq t\leq T\) reserve a common effective relatively
ample rational summand. The joint rational-family construction in
Proposition~\ref{prog:projective-good-model} puts these fixed pairs
in the polytopes required by Theorem~E of
\cite{FujinoAnalyticBCHM}.

Suppose a finite prefix has reached \(X_i\), and put
\(\mu_i=\lambda_{i-1}\), with \(\mu_0=T\). Its transformed
\(D_i+\mu_iN_i\) is nef. Define
\[
 \lambda_i=\min\{t\in[0,\mu_i]:D_i+tN_i
                         \text{ is nef over }S\}.
\]
The feasible set is a nonempty closed interval. Inductively
\(\lambda_i>2\epsilon\): all previous steps have nonpositive
comparison for \(D+2\epsilon N\), so its transform is still not
pseudoeffective. A first threshold at most \(2\epsilon\) would
make this transform nef by convexity between the current and
preceding thresholds, a contradiction.

There is a \(D_i\)-negative extremal ray \(R_i\) with
\((D_i+\lambda_iN_i)\cdot R_i=0\). To see attainment without a
generality assumption, let \(t<\lambda_i\) tend to \(\lambda_i\).
A ray negative for \(D_i+tN_i\) has positive \(N_i\)-degree,
because \(D_i+\mu_iN_i\) is nef. When \(t>\lambda_i/2\), it is
also negative for \(D_i+(\lambda_i/2)N_i\). The latter has fixed
ordinary klt big-boundary representatives on the charts, since
\(\lambda_i/2>\epsilon\). Normalize each of these boundaries by
\cite[Lemma~11.15]{FujinoAnalyticBCHM} as \(A_\ell+B'_\ell\),
with \(A_\ell\geq0\) relatively ample and the new pair klt.
The \(A_\ell\)-truncation of the cone theorem for
\(K+B'_\ell\) gives finitely many rays negative for the entire
adjoint in question. The finite cover therefore gives finitely
many global negative rays. One of these rays attains the wall,
as required.

Choose one such ray, even if the wall vanishes on other rays.
The rational-support argument for an individual negative ray in
Proposition~\ref{prog:relative-klt} gives a global rational nef
support annihilating exactly \(R_i\), and ordinary relative base
point freeness contracts it. For the real ordinary boundary one
may first choose a nearby rational klt boundary that is still
negative on \(R_i\). The contraction therefore has a rational
ordinary antiample adjoint. Its flip, when small, is constructed
from one global detector algebra as above.

If the contraction is of fibre type, stop. Otherwise the wall
\(w_i=D_i+\lambda_iN_i\) has degree zero on \(R_i\).
In the finite rational span of its actual line factors, this
zero-degree subspace is rational: its defining functional has
rational values on rational lines. Thus \(w_i\) is a real
combination of rational lines of zero ray degree.
Lemma~\ref{prog:integral-descent} descends those lines to the
contraction base. Their combination is nef over \(S\), by lifting
curves through the projective contraction. Pulling it to the
positive side shows that the transformed wall stays nef. The next
threshold therefore satisfies \(\lambda_{i+1}\leq\lambda_i\);
equality is permitted. No interval of ampleness between successive
thresholds is required.

If \(E_i\geq0\) is the ordinary \(D_i\)-comparison on a common
resolution, this actual wall descent gives, for every real \(t\),
\[
 p_i^*(D_i+tN_i)
 \sim_{\R}q_i^*(D_{i+1}+tN_{i+1})
                 +(1-t/\lambda_i)E_i.
\]
For \(t\leq\lambda_i\) the error is effective. This proves all
the nonpositivity and fixed-pair persistence assertions used in
the induction, including the lower bound on thresholds.

At every working model use the \emph{wall parameter} \(t=\lambda_i\),
rather than an interior ample parameter. Since
\(\lambda_i\leq\lambda_j\) for every earlier step, the displayed
comparisons show that \(X_i\) is a weak log canonical model of
the initial fixed chart pairs representing \(D+\lambda_iN\).
Its trace is nef by definition. Theorem~E counts all these marked
weak log canonical models, including the ones at repeated wall
parameters. The finite chart cover gives a finite global marked
list. Repetition would contradict the strict ordinary
\(D\)-discrepancy increase at an intervening step. Thus this
program with the specified \(N\) is finite. Non-pseudoeffectivity
excludes a nef endpoint for \(D\), so the last contraction is a
Mori fibre contraction. The comparison at \(t=1=\lambda_i\) is
crepant: such a step supplies no strict discrepancy improvement
for \(D+N\), while it is still strictly negative for \(D\).
\end{proof}
